\chapter{Kikuchi LDC — Defs}
%%% generated by the blueprint pipeline from TCSlib.KikuchiLDC.LDC.Defs
\section{Overview}

This module sets up the core objects used in the Alrabiah--Guruswami--Kothari--Manohar
near-cubic lower bound for $3$-query locally decodable codes via semirandom CSP
refutation: hypergraphs on a vertex set $\mathrm{Fin}\,n$ together with uniformity,
matching and degree notions; the normal form of a $(3,\delta,\varepsilon)$-locally
decodable code as a family of $3$-uniform matchings; the associated XOR polynomial and
its value over $\pm 1$ assignments; and the double-copy notation on $\mathrm{Fin}\,n
\times \mathrm{Fin}\,2$ used later in the Kikuchi matrix construction.

\section{Declarations}

\begin{definition}[Hypergraph on $\mathrm{Fin}\,n$]
\lean{Hypergraph}
\leanok
\statementsource{arXiv.2308.15403}{pdf-b372a20892ce-p004-b008}
A hypergraph on the vertex set $\mathrm{Fin}\,n$ is a finite collection of subsets of
$\mathrm{Fin}\,n$, modelled as $\mathtt{Finset}\,(\mathtt{Finset}\,(\mathrm{Fin}\,n))$.
Its elements are called edges (or constraints).
\end{definition}

\begin{definition}[$q$-uniform hypergraph]
\lean{Hypergraph.IsUniform}
\leanok
\statementsource{arXiv.2308.15403}{pdf-b372a20892ce-p004-b008}
\uses{Hypergraph}
A hypergraph $H$ on $\mathrm{Fin}\,n$ is $q$-uniform when every edge $C \in H$ has
exactly $q$ vertices, i.e.\ $\abs{C} = q$.
\end{definition}

\begin{definition}[Matching]
\lean{Hypergraph.IsMatching}
\leanok
\statementsource{arXiv.2308.15403}{pdf-b372a20892ce-p004-b008}
\uses{Hypergraph}
A hypergraph $H$ is a matching when its edges are pairwise disjoint: for all
$C, C' \in H$ with $C \neq C'$, the sets $C$ and $C'$ are disjoint.
\end{definition}

\begin{definition}[Degree of a set in a hypergraph]
\lean{Hypergraph.degree}
\leanok
\statementsource{arXiv.2308.15403}{
  pdf-b372a20892ce-p004-b008,
  pdf-b372a20892ce-p006-b002
}
\uses{Hypergraph}
The degree of a vertex set $Q \subseteq \mathrm{Fin}\,n$ in a hypergraph $H$ is the
number of edges of $H$ containing $Q$:
\[
  \deg_H(Q) \;=\; \abs{\{\, C \in H \;:\; Q \subseteq C \,\}}.
\]
\end{definition}

\begin{definition}[Pair-degree bound]
\lean{Hypergraph.pairDegreeBound}
\leanok
\statementsource{arXiv.2308.15403}{
  pdf-b372a20892ce-p006-b005,
  pdf-b372a20892ce-p006-b003
}
\uses{Hypergraph, Hypergraph.degree}
A hypergraph $H$ satisfies the pair-degree bound $d$ when every pair of distinct
vertices $u \neq v$ lies in at most $d$ edges, i.e.\ $\deg_H(\{u,v\}) \le d$.
\end{definition}

\begin{definition}[Sign values $\pm 1$]
\lean{IsPMOne}
\leanok
The predicate on an integer $x$ stating that $x$ is a sign value, i.e.\ $x = 1$ or
$x = -1$. Assignments used below take values in this set.
\end{definition}

\begin{definition}[Normal form locally decodable code]
\lean{NormalLDC}
\leanok
\statementsource{LDC_now}{pdf-4377a4b7ed99-p075-b002}
\statementsource{arXiv.2308.15403}{pdf-b372a20892ce-p005-b002}
\uses{Hypergraph, Hypergraph.IsMatching, Hypergraph.IsUniform}
The structural data of a $(3,\delta,\varepsilon)$-normally decodable code with message
length $k$ and block length $n$: for each message index $i : \mathrm{Fin}\,k$ a
hypergraph $H_i$ on $\mathrm{Fin}\,n$ which is $3$-uniform and a matching, together
with real parameters $\delta > 0$ and $\varepsilon > 0$ and the density requirement
$\delta n \le \abs{H_i}$ for every $i$.
\end{definition}

\begin{definition}[Combined hypergraph of an LDC]
\lean{NormalLDC.combined}
\leanok
\statementsource{arXiv.2308.15403}{pdf-b372a20892ce-p005-b008}
\uses{Hypergraph, NormalLDC}
For a normal form code $L$, the hypergraph obtained as the union of all its matchings,
$\bigcup_{i : \mathrm{Fin}\,k} H_i$, viewed as a hypergraph on $\mathrm{Fin}\,n$.
\end{definition}

\begin{definition}[Total number of constraints]
\lean{NormalLDC.totalConstraints}
\leanok
\statementsource{arXiv.2308.15403}{pdf-b372a20892ce-p005-b008}
\uses{NormalLDC}
For a normal form code $L$, the total number of constraints
\[
  m \;=\; \sum_{i : \mathrm{Fin}\,k} \abs{H_i}.
\]
\end{definition}

\begin{definition}[Monomial of an assignment on a set]
\lean{assignment_prod}
\leanok
For an assignment $x : \mathrm{Fin}\,n \to \mathbb{Z}$ and a set
$C \subseteq \mathrm{Fin}\,n$, the product
\[
  x_C \;=\; \prod_{v \in C} x_v .
\]
\end{definition}

\begin{definition}[XOR polynomial of a message]
\lean{NormalLDC.xorPoly}
\leanok
\statementsource{arXiv.2109.04415}{pdf-9c5df3bb4a04-p026-b001}
\statementsource{arXiv.2308.15403}{pdf-b372a20892ce-p005-b009}
\uses{NormalLDC, NormalLDC.totalConstraints, assignment_prod}
For a normal form code $L$, a message $b : \mathrm{Fin}\,k \to \mathbb{Z}$ and an
assignment $x : \mathrm{Fin}\,n \to \mathbb{Z}$, the real-valued XOR instance
\[
  \psi_b(x) \;=\; \frac{1}{m} \sum_{i : \mathrm{Fin}\,k} b_i
    \sum_{C \in H_i} x_C ,
\]
where $m$ is the total number of constraints of $L$.
\end{definition}

\begin{definition}[Value of the XOR instance]
\lean{NormalLDC.xorVal}
\leanok
\statementsource{arXiv.2109.04415}{
  pdf-9c5df3bb4a04-p026-b001,
  pdf-9c5df3bb4a04-p007-b001
}
\statementsource{arXiv.2308.15403}{pdf-b372a20892ce-p005-b009}
\uses{IsPMOne, NormalLDC, NormalLDC.xorPoly}
The value $\mathrm{val}(\psi_b)$ of the XOR instance associated with a message $b$,
defined as the supremum of $\psi_b(x)$ over all assignments
$x : \mathrm{Fin}\,n \to \mathbb{Z}$ taking values in $\{-1,1\}$.
\end{definition}

\begin{definition}[Double-copy ground set]
\lean{DoubleCopy}
\leanok
\statementsource{arXiv.2109.04415}{
  pdf-9c5df3bb4a04-p017-b003,
  pdf-9c5df3bb4a04-p029-b005
}
\statementsource{arXiv.2308.15403}{pdf-b372a20892ce-p007-b009}
The doubled vertex set $\mathrm{Fin}\,n \times \mathrm{Fin}\,2$, carrying two labelled
copies of each vertex of $\mathrm{Fin}\,n$.
\end{definition}

\begin{definition}[First copy of a vertex]
\lean{copy1}
\leanok
\statementsource{arXiv.2109.04415}{pdf-9c5df3bb4a04-p017-b003}
\statementsource{arXiv.2308.15403}{pdf-b372a20892ce-p007-b009}
\uses{DoubleCopy}
The map sending a vertex $u : \mathrm{Fin}\,n$ to its first copy
$u^{(1)} = (u,0)$ in the double-copy ground set.
\end{definition}

\begin{definition}[Second copy of a vertex]
\lean{copy2}
\leanok
\statementsource{arXiv.2109.04415}{pdf-9c5df3bb4a04-p017-b003}
\statementsource{arXiv.2308.15403}{pdf-b372a20892ce-p007-b009}
\uses{DoubleCopy}
The map sending a vertex $u : \mathrm{Fin}\,n$ to its second copy
$u^{(2)} = (u,1)$ in the double-copy ground set.
\end{definition}

\begin{definition}[First-copy embedding of a set]
\lean{liftCopy1}
\leanok
\statementsource{arXiv.2109.04415}{
  pdf-9c5df3bb4a04-p017-b003,
  pdf-9c5df3bb4a04-p029-b005
}
\statementsource{arXiv.2308.15403}{pdf-b372a20892ce-p007-b009}
\uses{DoubleCopy, copy1}
The image $C^{(1)}$ of a set $C \subseteq \mathrm{Fin}\,n$ under the injection
$u \mapsto u^{(1)}$, a subset of the double-copy ground set.
\end{definition}

\begin{definition}[Second-copy embedding of a set]
\lean{liftCopy2}
\leanok
\statementsource{arXiv.2109.04415}{
  pdf-9c5df3bb4a04-p017-b003,
  pdf-9c5df3bb4a04-p029-b005
}
\statementsource{arXiv.2308.15403}{pdf-b372a20892ce-p007-b009}
\uses{DoubleCopy, copy2}
The image $C^{(2)}$ of a set $C \subseteq \mathrm{Fin}\,n$ under the injection
$u \mapsto u^{(2)}$, a subset of the double-copy ground set.
\end{definition}
